\clearpage\begingroup\SourceLeaf{99b}\HistoricalHead{94}{L}{PHYSIK. BRIEFWECHSEL.}
\begin{center}[14.]\end{center}
\begin{center}\textsc{Gauss} an \textsc{Encke}.\quad Göttingen, 13. Oktober 1834.\end{center}
\noindent .\ .\ .\ .\ .\ .\ .\ .\ .\ .\ .\ .\ .\ .\ .\ .\ .\ .\ .\ .\ .
\par\noindent\hspace*{\parindent}Wenn Sie die 4 Suiten zeichnen, was auf dem carrirten Papier sehr schnell
geht, werden Sie bemerken, dass der Gang an allen vier Orten derselbe ist,
ohne eine erhebliche Störung. Inzwischen bemerkt man doch einige obwohl
sehr kleine Störungen, die umso interessanter sind, da bei denselben die
Gleichheit an allen vier Orten und mit Rücksicht auf den Meridianunterschied
unverkennbar ist. Besonders die kleine Vertiefung in der Nacht, die für
Göttingen etwa \(2^{\mathrm h}\,50'\), für Braunschweig \(2^{\mathrm h}\,52\frac{1}{2}'\), für Leipzig \(3^{\mathrm h}\,0'\) und Berlin
\(3^{\mathrm h}\,5'\) eingetreten ist; auch das gleichzeitige steile Herabgehen am ersten Vormittage
um \(11^{\mathrm h}\) ist frappant. Endlich hat meine Aufmerksamkeit erregt die
\clearpage\SourceLeaf{100}\HistoricalHead{95}{R}{GAUSS AN ENCKE.}
\noindent Ausschweifung an demselben Vormittag \(8^{\mathrm h}\,40'\) im M[agnetischen] O[bservatorium],
wobei ich gerade correspondirend beobachtete und bewundernswürdig
genau dieselbe Figur erhielt; die Figur in Berlin ist ziemlich ähnlich, während
in Leipzig sich nur eine schwache Andeutung zeigt.

Auf Einladung von \textsc{Weber} hat mein Sohn in mehreren verabredeten
Stunden am 1. u. 2. Oktober beobachtet; hier war das Glück günstiger. Besonders
am 1. Oktober Abends \(8\text{---}10^{\mathrm h}\) waren überaus grosse Anomalien, die
ich in der Sternwarte ganz übereinstimmend hatte. \textsc{Weber} wird wahrscheinlich
in diesen Tagen zurückkommen und die Leipziger Beobachtungen mitbringen,
auf die ich sehr neugierig bin.

Die Beobachtungen vom 23. 24. September zeigen aufs neue wie nothwendig
und wichtig es ist, in kleinen Zeitintervallen zu beobachten. Da jedoch
besonders an den Orten, wo nicht zahlreiche Gehilfen sind, solche Beobachtungen
44 Stunden, und 8mahl im Jahre zu machen, gar zu ermüdend
ist, so ist es wohl am besten, künftig eine andere Terminbestimmung zu
machen, und die bisherigen ganz fahren zu lassen. Ohnehin können die mit
\textsc{Gambeyschen} Apparaten und in grossen Zeitintervallen gemachten, für unsere
Zwecke als ziemlich unnütz betrachtet werden. Ich dächte 24 Stunden und
6mahl im Jahre wäre genug, und würden gewiss, wenn auch nicht jedesmal
doch oft genug, interessante Erscheinungen zu erwarten sein. Nach allen bisherigen
Erfahrungen kommen grosse Anomalien mehr bei Nacht vor; um also
solche nicht zu zerschneiden, würde am rathsamsten sein, von einem Mittag
bis zum folgenden zu beobachten; ausserdem hat dies den Vortheil, dass das
beschwerliche Beleuchtungs Arrangement nur einmahl getroffen zu werden
braucht. Sobald \textsc{Weber} zurück ist, werden wir das Nöthige deshalb concertiren,
und Ihnen anzeigen. Auch Prof. \textsc{Rosenberger} wünscht in Halle Theil
zu nehmen, und wird in Kurzem hier kommen, um die Beobachtungsart
einzulernen.
\par\noindent .\ .\ .\ .\ .\ .\ .\ .\ .\ .\ .\ .\ .\ .\ .\ .\ .\ .\ .\ .\ .
\par\noindent\hspace*{\parindent}Durch einen dichten Kasten hat meine 25 pfündige Nadel noch bedeutend
gewonnen. Ich werde bald scharfe Bestimmungen zur Ausmittlung des Temperatureinflusses
machen. Einen Multiplicator für die grosse Nadel von 250---300 Umwindungen lasse ich jetzt anfertigen. \mbox{.\ .\ .\ .\ .\ .\ .\ .\ .\ .\ .\ .\ .\ .\ .\ .\ .\ .\ \raisebox{.6ex}{.}}
\clearpage\endgroup
